\newcommand{\CeresSymbol}{\text{\mdseries\AstroSymbols ⚳}}
\SourceLeaf{212}%
\begin{center}
{\small\setlength{\tabcolsep}{5pt}
\begin{tabular}{r r|r|r|r}
\multicolumn{2}{c|}{Seeberg}&\shortstack{Berechnete\\Asc. rect. $\CeresSymbol$}&\shortstack{Unter-\\schied}&\shortstack{Berechnete\\Abweichung}\\\hline
1801 December&7&$178^0\ 33'\ 33''6$&$+3''0$&$11^0\ 47'\ 33''$ N\\
1802 Januar&11&$186\ 46\ 9.3$&$+19.3$&$11\ 15\ 41$\\
—&; 16&$187\ 28\ 3.1$&$+9.9$&$11\ 26\ 40$\\
—&22&$188\ 6\ 45.9$&$+20.1$&$11\ 45\ 18$\\
—&25&$188\ 21\ 6.5$&$+27.3$&$11\ 56\ 49$
\end{tabular}}
\end{center}
Da bei Bestimmung dieser Elemente keine andere Breiten gebraucht worden als die \textsc{Piazzi}'schen, die alle sehr nahe beim Knoten liegen: so kann man von der Neigung noch keine grosse Schärfe fordern. Inzwischen glaubt Dr.\ \textsc{Gauss} nach einem Überschlag der \textsc{Olbers}'schen Declinationen, dass sie kaum mehr als eine Minute fehlerhaft sein dürfte. Dies ist auch ein neuer Beweis für die vorzügliche Güte der \textsc{Piazzi}'schen Declinationen.

Diese VI\textsuperscript{ten} Elemente werden in der Folge noch immer bedeutende Verbesserungen bedürfen; diese werden immer sicherer ausfallen, je entfernter die Beobachtungen von einander liegen, als wenn Dr.\ \textsc{Gauss} jetzt schon den vorhandenen Beobachtungen seine Elemente möglichst genau anpassen wollte. Indessen ist Dr.\ \textsc{Gauss} ein so fertiger und unermüdeter Rechner, dass ihm die Bestimmung neuer Elemente gar keine Mühe zu machen scheint; denn ehe ich michs versah, und ehe ich ihm meine Beobachtung vom 26. Januar zugeschickt hatte, überraschte er mich schon mit einer VII\textsuperscript{ten} Bahn der \textit{Ceres}, indem er dazu statt meiner Beobachtung vom 16., die vom 25. Januar gebrauchte, und damit folgende (VII) Elemente herausbrachte.
\begin{center}
(VII)
\par\smallskip
{\small
$\text{Epoche 1801}\ \left\{\begin{array}{lr}
\text{Seeberger Meridian}\ .\ .&77^0\ 27'\ 36''5\\
\text{Palermer Meridian}\ .\ .&77\ 27\ 30.9
\end{array}\right.$
\par\smallskip
\begin{tabular}{l r}
tägliche mittlere tropische Bewegung $\ .$&$769''7924$\\
Logarithmus der halben grossen Axe $\ .$&$0.4424742$\\
Excentricität $\ .\ .\ .\ .\ .\ .\ .\ .\ .\ .$&$0.0814064$
\end{tabular}
\par\smallskip
$\left.\begin{array}{l}
\text{Sonnenferne}\\
\AscendingNode
\end{array}\right\}\ \text{1801 ruhend}\ \begin{array}{lr}
.\ .\ .\ .\ .&325^0\ 57'\ 15''\\
.\ .\ .\ .\ .&80\ 58\ 40
\end{array}$
\par\smallskip
\begin{tabular}{l r}
grösste Mittelpunktsgleichung $\ .\ .\ .$&$9\ 20\ 8.0$\\
Neigung $\ .\ .\ .\ .\ .\ .\ .\ .\ .\ .\ .$&$10\ 37\ 56.6$\\
{}[Handschrift] Epoche 1802 Palermo $\ .\ .$&$155\ 30\ 25.2$
\end{tabular}}
\end{center}
Hierbei sind abermals keine andere, als \textsc{Piazzi}'sche Breiten gebraucht worden; auch stellen diese Elemente die sämmtlichen Palermer Beobachtungen möglichst genau dar, und schliessen sich nun auch besser an meine Beobachtungen an, wie gegenwärtige Übersicht zeigt:
\begin{center}
{\small\setlength{\tabcolsep}{4pt}
\begin{tabular}{r r|r|r|r|r}
\multicolumn{2}{c|}{Seeberg}&\shortstack{Berechnete\\Ascens. rect.}&\shortstack{Unter-\\schied}&\shortstack{Berechnete\\Abweichung N.}&\shortstack{Unter-\\schied}\\\hline
1801 December&7&$178^0\ 33'\ 29''2$&$-1''4$&&\\
1802 Januar&11&$186\ 45\ 47.6$&$-2.3$&&\\
&16&$187\ 27\ 38.8$&$-14.4$&&\\
&22&$188\ 6\ 18.2$&$-7.6$&&\\
&25&$188\ 20\ 37.2$&$-2.0$&$11^0\ 56'\ 58''4$&$+35''4$\\
&26&$188\ 24\ 37.0$&$-12.5$&$12\quad2\quad8.4$&\\
&28&$188\ 31\ 25.7$&$-12.1$&$12\quad9\ 55.6$&$+14.3$\\
&29&$188\ 34\ 14.1$&$-4.0$&$12\ 14\ 34.1$&\\
&30&$188\ 36\ 38.4$&$-5.5$&$12\ 19\ 19.8$&$+19.1$\\
&31&$188\ 38\ 38.3$&$-7.1$&$12\ 24\ 15.3$&\\
Februar&3&$188\ 42\quad9.5$&$-3.5$&$12\ 39\ 56.6$&$+17.6$\\
&4&$188\ 42\ 30.7$&$-5.6$&$12\ 45\ 25.7$&\\
&5&$188\ 42\ 26.7$&$-3.4$&$12\ 51\quad2.6$&$+34.7$\\
&9&$188\ 38\quad1.6$&$-2.3$&$13\ 14\ 43.0$&$+22.0$
\end{tabular}}
\end{center}
\expandafter\def\csname HistoricalNoteText18\endcsname{Um auch hier nichts zu verabsäumen, was zur möglich-längsten Beobachtung dieses Planeten beitragen, und die möglich-grösste Genauigkeit verschaffen kann, werden wir jetzt, da es noch Zeit ist, alle die, in jenem Theil der scheinbaren Bahn der \textit{Ceres} zu liegen kommende Sterne vorher sehr genau zu bestimmen suchen, wo man nämlich den Planeten nicht mehr im Mittagskreis, sondern in seiner westlichen Quadratur, nur mit Aequatorial-Instrumenten oder parallactischen Fernröhren wird beobachten können. Auf jene Beobachtungen ist auch deswegen eine grössere Aufmerksamkeit zu verwenden, weil sie zur Bestimmung des \textit{Radius vector} die günstigsten sein werden.}
\SourceLeaf{213}%
Es wäre dem Dr.\ \textsc{Gauss} ein leichtes, obige geringe Abweichung der Rechnung von allen meinen bisher beobachteten geraden Aufsteigungen ganz und gar verschwinden zu machen. Allein da wir jetzt fast mit jeder Woche von der Verbesserung dieser Elemente einen grössern Grad von Genauigkeit hoffen können, und diese Elemente auf jeden Fall zur Auffindung der \textit{Ceres} auf lange Zeit hinlänglich genau sind: so glaubt Dr.\ \textsc{Gauss} selbst, dass es einstweilen am besten sein wird, die Vergleichung der Beobachtungen mit \textit{einerlei Elementen} erst noch eine Zeitlang fortzusetzen, und besonders erst noch eine etwas grössere Anzahl genauer Abweichungen abzuwarten. Selbst die geringe Abweichung, welche obige Tabelle darstellt, und die zu geringe Anzahl zuverlässiger Declinationen zeigen schon, dass es zur weiweitern Berichtigung der VII Elemente noch zu früh sei. Übrigens ist ja die Bahn der \textit{Ceres}, wegen der Perturbationen, welche sie von der Einwirkung der übrigen Planeten erleidet, keine vollkommene Ellipse. Diese Störungen müssen erst berechnet werden, ehe man die \textit{wahre} elliptische Bahn erhalten kann. —\quad —\quad —\quad Indessen, bei der jetzigen Lage der Dinge, wo die \textit{Ceres} bereits fast ein Viertel ihrer Bahn beschrieben hat, schmeichelt sich Dr.\ \textsc{Gauss}, dass der Unterschied zwischen diesen VII und den wahren Elementen der Bahn nicht sehr stark sein dürfte. Mit allen den Beobachtungen, die in diesem Jahre bis zum Unsichtbarwerden der \textit{Ceres} noch zu hoffen sind\HistoricalNote{18}, hofft er diese Elemente noch zu einem solchen Grade von Genauigkeit bringen zu können, dass man, wenn man 44 Jahre damit rückwärts bis zu \textsc{Mayer}'s Beobachtungen rechnet, nicht um einen Grad in der Länge, und nur wenige Minuten in der Breite, soll fehlen können. Sollte sich die \textit{Ceres} in keinem älteren Sternverzeichniss, so wie \textit{Uranus}, finden, so wird man freilich ihre mittlere Bewegung nur stufenweise genauer kennen lernen, und man wird dann noch lange Zeit hindurch jedes Jahr etwas hinzuzusetzen und zu verbessern haben. —\quad —\quad —

Um unsern astronomischen Lesern auch künftig die Aufsuchung der \textit{Ceres} zu erleichtern, so theilen wir ihnen hier abermals eine kleine \textit{Ephemeride} ihres Laufes für die nächsten zwei Monate mit, welche Dr.\ \textsc{Gauss} nach seinen letzten und VII Elementen berechnet hat. Mehr als eine Minute werden diese Stellungen hoffentlich im April vom Himmel nicht abweichen.
\begin{center}
Stellungen der \textit{Ceres} im März und April 1802 für Mitternacht, mittlere Seeberger Zeit.
\par\smallskip
{\small\setlength{\tabcolsep}{5pt}
\begin{tabular}{r r|r|r|r}
\multicolumn{2}{c|}{1802}&\shortstack{AR. $\CeresSymbol$\\in Graden}&Abweich.&AR. in Zeit\\\hline
März&1&$186^0\ 41'$&$15^0\ 30'$ N&$12^{\mathrm u}\ 26^{\mathrm m}\ 45^{\mathrm s}$\\
&4&$186\ 11$&$15\ 50$&$12\ 24\ 45$\\
&7&$185\ 39$&$16\ 10$&$12\ 22\ 36$\\
&10&$185\quad5$&$16\ 29$&$12\ 20\ 18$\\
&13&$184\ 28$&$16\ 47$&$12\ 17\ 53$\\
&16&$183\ 51$&$17\quad4$&$12\ 15\ 24$\\
&19&$183\ 13$&$17\ 19$&$12\ 12\ 50$\\
&22&$182\ 34$&$17\ 33$&$12\ 10\ 15$\\
&25&$181\ 55$&$17\ 44$&$12\quad7\ 40$
\end{tabular}
\par\smallskip
\begin{tabular}{r r|r|r|r}
\multicolumn{2}{c|}{1802}&\shortstack{AR. $\CeresSymbol$\\in Graden}&Abweich.&AR. in Zeit\\\hline
März&28&$181^0\ 17'$&$17^0\ 54'$ N&$12^{\mathrm u}\quad5^{\mathrm m}\quad7^{\mathrm s}$\\
&31&$180\ 39$&$18\quad1$&$12\quad2\ 37$\\
April&3&$180\quad3$&$18\quad6$&$12\quad0\ 12$\\
&6&$179\ 29$&$18\quad9$&$11\ 57\ 54$\\
&9&$178\ 56$&$18\ 10$&$11\ 55\ 45$\\
&12&$178\ 26$&$18\quad9$&$11\ 53\ 46$\\
&15&$177\ 59$&$18\quad5$&$11\ 51\ 57$\\
&18&$177\ 35$&$17\ 59\frac12$&$11\ 50\ 20$
\end{tabular}}
\end{center}
Diesen Ephemeriden fügt Dr.\ \textsc{Gauss} noch folgende daraus abgeleitete wichtige Bemerkungen bei:

„Der Gegenschein des Planeten fällt den 17. März Nachmittags. Um diese Zeit wird er am vortheilhaftesten zu beobachten sein. Kurz vor der Opposition hat er den kleinsten Abstand von der Erde
\SourceLeaf{214}%
erreicht $=1.6025$. Um dieselbe Zeit (etwas früher) fällt seine grösste nördliche geocentrische Breite $=17^0\ 9'$, und etwas später seine schnellste rückläufige Bewegung, täglich etwa 13 Min.\ in AR. Die nördliche Abweichung nimmt noch bis Anfang April zu; ungefähr den 9. April wird der Lauf wieder südlich.“ —\quad —\quad —
\begin{flushright}[\textsc{von Zach}.]\end{flushright}

\bigskip\Needspace{10\baselineskip}
\pdfbookmark[1]{Fortgesetzte Nachrichten. Ceres Ferdinandea. 1802 April}{ceres180204}
\begin{center}
{\small FORTGESETZTE NACHRICHTEN ÜBER DEN NEUEN HAUPTPLANETEN.\par}
CERES FERDINANDEA.\par
\medskip\hrule\medskip
Monatliche Correspondenz Band V. S.\ 379—404. 1802 April.\par
\medskip\hrule
\end{center}\medskip
—\quad —\quad —\quad Die \textsc{Gauss}'sche VII. elliptische Bahn des \textit{neuen Planeten} [S.\ 207 d.\ B.] stimmt fortdauernd genau mit unsern fortgesetzten Beobachtungen dieses planetarischen Weltkörpers, wie nachfolgende Berechnung und Vergleichung des Dr.\ \textsc{Gauss} zur Genüge beweist.
\begin{center}
{\small\setlength{\tabcolsep}{5pt}
\begin{tabular}{r r|r|r|r|r}
\multicolumn{2}{c|}{\shortstack{Seeberg\\1802}}&\shortstack{Berechnete\\Ascensio\\recta $\CeresSymbol$}&\shortstack{Unter-\\schied}&\shortstack{Berechnete\\Abweichung\\nördlich}&\shortstack{Unter-\\schied}\\\hline
Februar&19&$187^0\ 58'\ 23''8$&$-4''1$&$14^0\ 20'\ 24''0$&$+21''1$\\
—&26&$187\quad7\ 34.4$&$+1.0$&$15\quad9\ 20.1$&$+25.8$\\
—&27&$186\ 58\ 53.2$&$+5.8$&$15\ 16\ 20.2$&$+26.3$\\
—&28&$186\ 49\ 51.9$&$+3.7$&$15\ 23\ 19.1$&$+18.1$\\
März&1&$186\ 40\ 31.4$&$+3.5$&$15\ 30\ 16.0$&$+34.9$\\
—&2&$186\ 30\ 52.2$&$+3.3$&$15\ 37\ 10.6$&$\ .\ .\ .$\\
—&3&$186\ 20\ 55.1$&$+2.6$&$15\ 44\quad2.0$&$+19.7$
\end{tabular}}
\end{center}
Die geraden Aufsteigungen stimmen also noch immer vortrefflich, und der Fehler in der Abweichung scheint sich wenig oder gar nicht geändert zu haben. Mit Weglassung der vom 1. März wäre er im Mittel aus fünf Beobachtungen $+22''$. —\quad —\quad —

Der Oberamtmann Dr.\ \textsc{Schröter} fährt fort, den Durchmesser der \textit{Ceres} anhaltend zu beobachten. Folgendes sind einige seiner Beobachtungen, welche Dr.\ \textsc{Gauss} berechnet hat, und uns mitzutheilen so gütig war.
\begin{center}
{\small\setlength{\tabcolsep}{4pt}
\begin{tabular}{r l r l r}
\multicolumn{3}{c}{Scheinbarer Durchmesser}&\multicolumn{2}{c}{\shortstack{Berechneter Durchmesser\\in der Entfernung 1.}}\\[5pt]
1802 Januar 10&$\ .\ .\ .\ .$&$2''500$&$\ .\ .\ .\ .\ .\ .$&$5''16$\\
25&$\ .\ .\ .\ .$&$2.514$&$\ .\ .\ .\ .\ \raisebox{.5ex}{.}\ .$&$4.77$\\
26&$\ .\ .\ .\ .$&$2.687$&$\ .\ .\ .\ .\ .\ .$&$5.07$\\
28&$\ .\ .\ .\ .$&$2.774$&$\ .\ .\ .\ .\ .\ .$&$5.18$\\
31&$\ .\ .\ .\ .$&$2.930$&$\ .\ .\ .\ .\ .\ .$&$5.38$\\
Februar 5&$\ .\ .\ .\ .$&$3.120$&$\ .\ .\ .\ .\ .\ .$&$5.59$\\
10&$\ .\ .\ .\ .$&$3.342$&$\ .\ .\ .\ .\ .\ .$&$5.90$
\end{tabular}}
\end{center}
\SourceLeaf{215}%

Das Mittel wäre $5''29=0.308$ von der Erde, oder 529 geographische Meilen, mit Einschluss des. atmosphärischen Nebels. Nach \textsc{Schröter}'s Messung der scharf begrenzten Scheibe wäre der Durchmesser derselben nur $3''44$ in der Distanz $=1$, also fast genau $\tfrac15$ vom Durchmesser der Erde, oder beträchtlich kleiner, als unser Mond. Kein Wunder also, dass ein solches Planetchen uns so lange verborgen geblieben ist!

Da die \textsc{Gauss}'sche VII. Ellipse der \textit{Ceres}-Bahn die Beobachtungen dieses Planeten noch immerfort ziemlich genau darstellt, so theilen wir unsern astronomischen Lesern noch bei guter Zeit eine fortgesetzte \textit{Ephemeride} ihres Laufes für die künftigen drei Monate mit, deren Berechnung Dr.\ \textsc{Gauss} gütigst übernommen hat. Es ist zwar nicht glaubbar, dass \textit{Ceres} bis dahin sichtbar bleiben sollte; denn wir bemerken schon gegenwärtig, den 27. und 28. März, eine solche schnelle Lichtveränderung bei diesem Planeten, dass wir sie ganz andern Ursachen, als dem Zustande unseres Dunstkreises, oder seiner Entfernung von der Erde zuzuschreiben geneigt sind. Indessen, damit diejenigen Astronomen, welche mit Aequatorial- und parallactischen Instrumenten versehen sind, nichts unversucht lassen mögen, so werden sie sich die immer schwieriger werdende Aufsuchung der Ceres dadurch nicht wenig erleichtern.
\begin{center}
Stellungen der \textit{Ceres} im April, Mai und Juni 1802; für Mitternacht, mittlere Seeberger Zeit.\par\medskip
{\fontsize{9}{11}\selectfont\setlength{\tabcolsep}{2pt}
\begin{tabular}{r r|r|r|r@{\qquad}r r|r|r|r}
\multicolumn{2}{c|}{}&\shortstack{AR.\ $\CeresSymbol$\\in Graden}&\shortstack{Abweich.\\nördl.}&\shortstack{AR.\ $\CeresSymbol$\\in Zeit}&\multicolumn{2}{c|}{}&\shortstack{AR.\ $\CeresSymbol$\\in Graden}&\shortstack{Abweich.\\nördl.}&\shortstack{AR.\ $\CeresSymbol$\\in Zeit}\\\hline
April&21&$177^0\ 14'$&$17^0\ 52'$&$11^{\mathrm u}\ 48^{\mathrm m}\ 55^{\mathrm s}$&Mai&27&$177^0\ 14'$&$14^0\ 13'$&$11^{\mathrm u}\ 48^{\mathrm m}\ 57^{\mathrm s}$\\
&24&$176\ 56$&$17\ 42$&$11\ 47\ 43$&&30&$177\ 34$&$13\ 47$&$11\ 50\ 16$\\
&27&$176\ 41$&$17\ 30$&$11\ 46\ 44$&Juni&2&$177\ 57$&$13\ 21$&$11\ 51\ 47$\\
&30&$176\ 29$&$17\ 17$&$11\ 45\ 58$&&5&$178\ 22$&$12\ 53$&$11\ 53\ 27$\\
Mai&3&$176\ 21$&$17\quad2$&$11\ 45\ 25$&&8&$178\ 49$&$12\ 25$&$11\ 55\ 17$\\
&6&$176\ 17$&$16\ 45$&$11\ 45\quad6$&&11&$179\ 19$&$11\ 56$&$11\ 57\ 16$\\
&9&$176\ 15$&$16\ 27$&$11\ 45\quad2$&&14&$179\ 51$&$11\ 27$&$11\ 59\ 24$\\
&12&$176\ 17$&$16\quad8$&$11\ 45\ 10$&&17&$180\ 25$&$10\ 57$&$12\quad1\ 40$\\
&15&$176\ 23$&$15\ 47$&$11\ 45\ 27$&&20&$181\quad1$&$10\ 27$&$12\quad4\quad4$\\
&18&$176\ 31$&$15\ 25$&$11\ 46\quad5$&&23&$181\ 39$&$9\ 56$&$12\quad6\ 35$\\
&21&$176\ 43$&$15\quad2$&$11\ 46\ 50$&&26&$182\ 18$&$9\ 25$&$12\quad9\ 13$\\
&24&$176\ 57$&$14\ 38$&$11\ 47\ 48$&&29&$183\quad0$&$8\ 54$&$12\ 11\ 59$
\end{tabular}}
\end{center}

\textit{Ceres} kommt zum Stillstande, und wird rechtläufig den 2. Mai in Länge $5^{\mathrm z}\ 19^0\ 45'$, und den 9. Mai in AR $176^0\ 15'$. Die Grösse des Rückganges beträgt diesmal in Länge $13^0\ 10'$, und dauert 92 Tage, in AR $12^0\ 27'$, und dauert 93 Tage. Das Verhältniss der Lichtstärke des Planeten, mit Vernachlässigung der Grösse der Phase, hat Dr.\ \textsc{Gauss} folgendermaassen berechnet:
\begin{center}
{\small\setlength{\tabcolsep}{5pt}
\begin{tabular}{r l r|r|r|r}
\multicolumn{3}{c|}{}&\multicolumn{2}{c|}{Abstand von der}&\shortstack{Licht-\\stärke}\\
\multicolumn{3}{c|}{}&$\odot$&$\EarthSymbol$&\\\hline
1801&December&7&$2.5453$&$2.4892$&$0.4227$\\
1802&März&16&$2.5706$&$1.6024$&$1.0000$\\
&Mai&6&$2.5961$&$1.9039$&$0.6945$\\
&Juni&29&$2.6305$&$2.5735$&$0.3703$
\end{tabular}}
\end{center}
\begin{flushright}[\textsc{von Zach}.]\end{flushright}
\SourceLeaf{216}%
\bigskip\Needspace{10\baselineskip}
\pdfbookmark[1]{Fortgesetzte Nachrichten. Ceres Ferdinandea. 1802 Mai}{ceres180205}
\begin{center}
{\small FORTGESETZTE NACHRICHTEN ÜBER DEN NEUEN HAUPT-PLANETEN.\par}
CERES FERDINANDEA.\par
\medskip\hrule\medskip
Monatliche Correspondenz. Band V. S.\ 462—481. 1802 Mai.\par
\medskip\hrule
\end{center}\medskip
—\quad —\quad —\quad Bis zum 3. März haben wir im vorigen Hefte S.\ 389 [S.\ 209 d.\ B.] die Übereinstimmung unserer Beobachtungen mit den \textsc{Gauss}'schen Elementen angeführt. Seitdem hat Dr.\ \textsc{Gauss} diese Vergleichung, wie sie hier folgt, fortgesetzt.
\begin{center}
{\small\setlength{\tabcolsep}{5pt}
\begin{tabular}{r r|r|r|r|r}
\multicolumn{2}{c|}{\shortstack{Seeberg\\1802}}&\shortstack{Berechnete\\AR der $\CeresSymbol$}&\shortstack{Unter-\\schied}&\shortstack{Berechnete\\Abweich.\ der $\CeresSymbol$}&\shortstack{Unter-\\schied}\\\hline
März&6&$185^0\ 49'\ 23''3$&$+4''6$&$16^0\quad4'\ 12''2$&$+23''0$\\
—&7&$185\ 38\ 22.0$&$+6.1$&$16\ 10\ 45.6$&$+29.7$\\
—&10&$185\quad3\ 57.2$&$+8.3$&$16\ 29\ 48.4$&$+29.5$\\
—&11&$184\ 52\quad5.2$&$+9.4$&$16\ 35\ 55.2$&$\ .\ .\ .\ .\ .$\\
—&15&$184\quad3\quad2.9$&$+10.9$&$16\ 58\ 58.2$&$+27.5$\\
—&16&$183\ 50\ 28.4$&$+9.4$&$17\quad4\ 20.6$&$+21.7$\\
—&17&$183\ 37\ 48.5$&$+13.7$&$17\quad9\ 32.6$&$+26.2$\\
—&18&$183\ 25\quad3.2$&$+12.5$&$17\ 14\ 34.4$&$+29.1$\\
—&19&$183\ 12\ 14.4$&$+15.8$&$17\ 19\ 25.1$&$+34.3$
\end{tabular}}
\end{center}

Auch zwei \textit{Greenwicher} Beobachtungen verglich Dr.\ \textsc{Gauss} mit seinen Elementen, und die geraden Aufsteigungen stimmen auf das allergenaueste mit den unsrigen, wenn man obige unbezweifelte Voraussetzung macht, dass bei der vom 19. Februar durch ein Versehen die Minute verschrieben sei.
\begin{center}
{\small\setlength{\tabcolsep}{2.8pt}
\begin{tabular}{l r|r|r|r|r|r|r}
\multicolumn{2}{c|}{}&&\multicolumn{2}{c|}{Unterschied bei}&&\multicolumn{2}{c}{Unterschied bei}\\
\multicolumn{2}{c|}{\shortstack{Greenwich\\1802}}&\shortstack{Berechnete\\AR der $\CeresSymbol$}&\shortstack{Dr.\ \textsc{Mas-}\\\textsc{kelyne}}&\textsc{v. Zach}&\shortstack{Berechnete\\Abweichung\\der $\CeresSymbol$}&\shortstack{Dr.\ \textsc{Mas-}\\\textsc{kelyne}}&\textsc{v. Zach}\\\hline
Februar&19&$187^0\ 58'\ 13''3$&$-4''3$&$-4''1$&$14^0\ 20'\ 36''2$&$+35''3$&$+21''1$\\
März&6&$185\ 49\quad3.8$&$+5.0$&$+4.6$&$16\quad4\ 24.0$&$+34.5$&$+23.0$
\end{tabular}}
\end{center}

Diese Vergleichung zeigt nun offenbar an, dass sich die Beobachtungen des Planeten gegen Ende März zunehmend von der \textsc{Gauss}'schen Ellipse entfernen werden. Allein sollten sie auch, was doch nicht wahrscheinlich ist, im Juni um 5 Minuten fehlen, so wird dies doch \textit{gar nichts} auf sich haben, denn der Fehler kann sich nur allmählich anhäufen, und man wird ihn also, wenn man immer die vorhergehenden Beobachtungen mit den Elementen oder mit der Ephemeride vergleicht, allemal mit. hinreichender Schärfe voraus wissen. Dr.\ \textsc{Gauss} schreibt uns auch, dass er bereits solche Vorbereitungen gemacht habe, dass er die Elemente gleich wieder so ändern könne, dass sie sich an die neuesten Beobachtungen auf das genaueste anschliessen, sobald der Fehler sich auf eine Minute anhäufen sollte, und man kann alsdann wieder auf eine geraume Zeit eine sehr genaue Übereinstimmung erwarten. Sobald aber die Beobachtungen für dies Jahr geendigt und beisammen sein werden, gedenkt er die wahre Ellipse mit allen ihren Perturbations-Gleichungen so genau als möglich zu bestimmen.
\begin{flushright}[\textsc{von Zach}.]\end{flushright}
\SourceLeaf{217}%
\bigskip\Needspace{12\baselineskip}
\pdfbookmark[1]{Fortgesetzte Nachrichten über einen neuen von Dr. Olbers entdeckten Planeten. 1802 Juni}{pallas180206}
\fancyhead[L]{\small PALLAS}
\begin{center}
{\small FORTGESETZTE NACHRICHTEN ÜBER EINEN NEUEN VON DR. OLBERS IN BREMEN ENTDECKTEN ZWISCHEN MARS UND JUPITER SICH BEWEGENDEN HÖCHST MERKWÜRDIGEN PLANETEN UNSERES SONNEN-SYSTEMS.\par}
\medskip\hrule\medskip
Monatliche Correspondenz. Band V. S.\ 591—606. 1802 Juni.\par
\medskip\hrule
\end{center}\medskip
Als wir im vorhergehenden Hefte die erste Anzeige der merkwürdigen Entdeckung des neuen \textsc{Olbers}'schen Gestirns machten: so berechtigten uns die damals vorhandenen Beobachtungen, Berechnungen und gesammelten Erfahrungen noch nicht, diesen sonderbaren Wandelstern als einen, unserem Sonnen-Systeme angehörigen Planeten zu verkündigen; obgleich schon Vermuthungen, welche wir auch geäussert haben, vorhanden waren, dass dieses seiner Bewegung und seinem äussern Ansehen nach sich so sehr auszeichnende Gestirn nicht in die Classe der gewöhnlichen sogenannten Cometen zu setzen sei. So gross war indessen schon die Vermuthung eines \textit{Planetismus} gleich bei dem ersten Anblicke dieses Weltkörpers, dass ihr Entdecker es sogleich wagte, dessen Bahn, nicht etwa erst, wie es sonst zu geschehen pflegt, in einer Parabel, sondern in einem mit der Sonne concentrischen Kreise zu versuchen. Allein es liessen sich, wie wir schon S.\ 486 des vorigen Heftes angezeigt haben, die vorhandenen Beobachtungen durchaus in keiner Kreisbahn darstellen, und so gering auch ihre Anzahl war, so machte es die starke Bewegung dieses Gestirns in der Breite doch schon möglich, den Kreis bei künftigen Untersuchungen dieser Bahn auf immer auszuschliessen.

Dasselbe versuchte Dr.\ \textsc{Gauss}; auch er fand Dr.\ \textsc{Olbers} Erfahrung bestätigt, dass nämlich die Bewegung der \textit{Pallas} immer schneller war, als sie den \textsc{Kepler}'schen Gesetzen zu Folge in einem Kreise sein konnte. Sobald aber Dr.\ \textsc{Gauss} des Dr.\ \textsc{Olbers} Beobachtung dieses Gestirns vom 17. April erhalten hatte, unternahm er die Bestimmung des Kegelschnittes unabhängig von aller Hypothese, nach seiner eigenen Methode, deren Vorzüglichkeit schon \textit{Ceres} auf eine so glänzende Art bestätigt hat. Es liess sich zwar voraussehen, dass die kreismikrometrischen Beobachtungen zu dieser delicaten Rechnung nicht genau, und ihre Dauer nicht lang genug sei, zumal da \textit{Pallas} im Anfang Aprils in einer ungünstigen Lage war, indem die Tangente der geocentrischen Bewegung durch den Sonnenort ging. Inzwischen fand Dr.\ \textsc{Gauss} eine Ellipse, in ihren Hauptdimensionen von der, die nachher bestimmt wurde, nicht sehr verschieden, und das Resultat stimmte im allgemeinen dahin überein, dass: Pallas ein Planet zwischen Mars und Jupiter sei, dessen Bahn an einer Stelle der Ceres-Bahn sehr nahe kommt, und bei einer sehr starken Neigung eine beträchtliche Excentricität hat.

Die Excentricität war hiebei $0.3$; also stärker als bei allen andern Planeten, aber doch nicht gross genug, um der \textit{Pallas} den Namen eines Planeten verweigern zu können. Denn das Verhältniss der grossen zur kleinen Axe ist wie $1:\BareRoot{(1-(0.3)^2)}$; also ungefähr wie $1:0.95$; folglich die Ellipse noch nicht so platt, als die Jupiter-Scheibe.

Da die hiezu gebrauchten Beobachtungen natürlich bei weitem nicht so genau als Meridian-Beobachtungen sind, so setzte Dr.\ \textsc{Gauss} ein gegründetes Misstrauen in diese Elemente, und hielt ihre Bekanntmachung bis auf künftige Prüfung zurück. Erst wie er unsere \textit{Seeberger} Meridian-Beobachtungen dieses Gestirns vom 15. 18. 19. April, und nachher drei Pariser Beobachtungen vom 10. 12. 13. April erhielt, entschloss er sich, in Ermangelung einer längeren Reihe von genauen Beobachtungen, diese vor%
\SourceLeaf{218}%
läufigen Elemente so zu verbessern, dass sie sich an sechs \textit{Seeberger} Beobachtungen vom 4. 5. 7. 15. 18. 19. April, und an die drei \textit{Pariser} anschlössen, und so fand er sofort beim ersten Versuch folgende
\expandafter\def\csname HistoricalNoteText19\endcsname{Eben die Voraussetzung, welche Dr.\ \textsc{Gauss} auch bei der \textit{Ceres} zum Grunde gelegt hatte, wo er das Aphel und den $\AscendingNode$ nicht in Beziehung auf unser Aequinoctium, sondern auf die Fixsterne als ruhend ansah.}
\Needspace{20\baselineskip}
\begin{center}
\textit{Elemente der Pallas} (I).\par\medskip
[Handschrift:] Heliocentrischer Bogen $4^0\ 11'$\par\medskip
Bahn einer rechtläufigen Ellipse.\par\medskip
{\small\setlength{\tabcolsep}{2pt}
\begin{tabular}{l r}
Epoche 1802 März 31 Mittag in Seeberg $\ .\ .\ .\ .\ .\ .\ .$&$166^0\quad1'\ 37''2$\\
Tägliche mittlere tropische Bewegung $\ .\ .\ .\ .\ .\ .\ .$&$800''770$\\
Umlaufszeit $\ .\ .\ .\ .\ .\ .\ .\ .\ .\ .\ .\ .\ .\ .\ .\ .\ .$&$1618\tfrac12$ Tage\\
Logarithmus der halben grossen Axe $\ .\ .\ .\ .\ .\ .\ .$&$0.4310494$
\end{tabular}
\par\smallskip
$\left.\begin{array}{l}\text{Sonnenferne}\\\text{Knoten}\end{array}\right\}\quad
\text{\shortstack{zur Zeit der Epoche\\siderisch ruhend\footnotemark}}\left\{\begin{array}{lr}.&304^0\ 36'\ 30''\\.&172^0\quad9'\ 58''\end{array}\right.$
\footnotetext{\csname HistoricalNoteText19\endcsname}
\par\smallskip
\begin{tabular}{l r}
Excentricität $\ .\ .\ .\ .\ .\ .\ .\ .\ .\ .\ .\ .\ .\ .\ .\ .$&$0.215708$\\
Neigung $\ .\ .\ .\ .\ .\ .\ .\ .\ .\ .\ .\ .\ .\ .\ .\ .\ .$&$33^0\ 39'\ 16''6$
\end{tabular}}
\end{center}
Mit diesen Elementen verglich nun Dr.\ \textsc{Gauss} sehr scharf sowol die \textit{Seeberger} Beobachtungen, welche er damals hatte, als auch die übrigen von uns erhaltenen, zugleich mit den \textit{Pariser} Beobachtungen und einer \textsc{Olbers}'schen, welche folgende Übereinstimmung gaben.
\begin{center}
{\small\setlength{\tabcolsep}{3pt}
\begin{tabular}{r r|r|r|r|r|l}
\multicolumn{2}{c|}{}&&&\multicolumn{2}{c|}{Unterschiede}&\\
\multicolumn{2}{c|}{1802}&\shortstack{Berechnete\\AR.\ der Pallas}&\shortstack{Berechnete\\Abweichung\\der Pallas}&in AR.&in Decl.&Beobachter\\\hline
April&4&$183^0\ 44'\quad6''3$&$13^0\ 54'\ 51''0$&$-0''3$&$-2''0$&von Zach\\
&5&$183\ 34\ 25.7$&$14\ 13\ 17.2$&$+2.0$&$-5.7$&—\\
&7&$183\ 15\ 40.0$&$14\ 48\ 55.5$&$+1.5$&$-6.6$&—\\
&8&$183\quad6\ 35.2$&$15\quad6\quad9.1$&$-2.6$&$-0.9$&—\\
&10&$182\ 48\ 56.5$&$15\ 39\ 40.9$&$-1.5$&$-1.1$&Dr.\ Burckhardt\\
&12&$182\ 32\ 34.3$&$16\ 11\quad4.2$&$+2.8$&$+6.2$&Méchain\\
&13&$182\ 24\ 47.2$&$16\ 26\quad6.6$&$+2.8$&$+0.6$&—\\
&15&$182\ 10\ 13.2$&$16\ 54\ 33.0$&$-3.3$&$+2.2$&von Zach\\
&18&$181\ 50\ 26.2$&$17\ 34\ 26.5$&$-4.4$&$\ .\ .\ .$&—\\
&19&$181\ 44\ 28.0$&$17\ 46\ 51.7$&$+2.7$&$-2.5$&—\\
&24&$181\ 19\ 39.4$&$18\ 42\ 27.8$&$+1.6$&$\ .\ .\ .$&—\\
&25&$181\ 15\ 44.1$&$18\ 52\ 18.3$&$+11.9$&$\ .\ .\ .$&—\\
&26&$181\ 12\ 10.4$&$19\quad1\ 43.5$&$+8.6$&$-5.9$&—\\
&27&$181\quad8\ 57.9$&$19\ 10\ 44.0$&$+12.2$&$-2.8$&—\\
&29&$181\quad3\ 39.6$&$19\ 27\ 31.2$&$+23.0$&$-12.5$&—\\
&30&$181\quad1\ 33.9$&$19\ 35\ 18.2$&$+24.0$&$-15.3$&—\\
Mai&1&$180\ 59\ 48.1$&$19\ 42\ 42.5$&$+25.6$&$+1.7$&—\\
&1&$180\ 59\ 38.0$&$19\ 43\ 38.0$&$+20.0$&$+7.0$&Dr.\ Olbers
\end{tabular}}
\end{center}
Die Unterschiede mit verkehrten Zeichen der Rechnung zugesetzt geben die Beobachtung.

Dr.\ \textsc{Gauss} behält auch hier, wie bei der \textit{Ceres}, seine alte Gewohnheit bei (welche Nachahmung verdient), die geraden Aufsteigungen und Abweichungen aus den Elementen selbst abzuleiten, und mit den beobachteten zu vergleichen, und nicht aus den beobachteten AR und Decl.\ \textit{beobachtete Längen und Breiten herzuleiten}. Der Grund dieses Verfahrens ist so einleuchtend, dass man sich wundern muss, dass
\SourceLeaf{219}%
solches nicht schon längst allgemein eingeführt und befolgt worden ist. In Ansehung der Mühe ist es gleich viel, ob man mit der Beobachtung der Rechnung einen Schritt entgegen geht, oder die Rechnung ganz bis zur Beobachtung hinführt. Nach einem dem Dr.\ \textsc{Gauss} eigenen Verfahren wird diese Parallaxen-Rechnung vielmehr noch etwas einfachen. Man gewinnt aber dadurch hauptsächlich den Vortheil, das Gute von dem etwa weniger Zuverlässigen zu scheiden; sonst kann es leicht kommen, dass eine gute gerade Aufsteigung durch eine unrichtige Declination verdorben wird, und so nicht blos eine fehlerhafte Breite, sondern auch eine schlechte Länge herauskommt. Da bei dem obigen Vergleich der Beobachtungen mit den Elementen der Bahn die Unterschiede gegen Ende immer zunehmend werden, so liess sich voraussehen, dass bei der kurzen Dauer der Beobachtungen ein Unterschied von $25''$ noch einen sehr starken Einfluss auf die Bestimmung dieser Elemente haben müsse. Eine fortgesetzte Untersuchung, die Dr.\ \textsc{Gauss} nach dieser Vergleichung sogleich wieder vornahm, hatte ihm auch gezeigt, dass die 15tägigen Beobachtungen, auf welche sich seine vorigen Elemente gründeten, zwar nicht hinlänglich waren, eine schon genaue Bahn-Bestimmung zu liefern, aber doch die Natur des Gestirns schon ziemlich kenntlich anzeigten. Er hat nun sofort aus der ganzen Reihe unserer \textit{Seeberger} Beobachtungen neue Elemente bestimmt, die sich, so gut sichs thun liess, an diese Beobachtungen vom 4. April bis 1. Mai anschliessen, und da der auf der Bahn durchlaufene Bogen bereits $7^0\ 26'$ beträgt, so dürfte man zum wenigsten über den Planetismus des \textsc{Olbers}'schen Gestirns wol keinen Zweifel mehr hegen. Die neue Ellipse dieses planetarischen Weltkörpers ist demnach folgende:
\begin{center}
\textit{Elemente der Pallas aus 27tägigen Seeberger Beobachtungen.}\par\medskip
(II.)\par\medskip
{\small\setlength{\tabcolsep}{2pt}
\begin{tabular}{l r}
Epoche 1802 März 31 Mittag in Seeberg $\ .\ .\ .\ .\ .\ .\ .$&$161^0\ 12'\ 43''2$\\
Tägliche mittl.\ tropische Bewegung $\ .\ .\ .\ .\ .\ .\ .\ .$&$757''166$\\
Logarithmus der halben grossen Axe $\ .\ .\ .\ .\ .\ .\ .\ .$&$0.4472636$
\end{tabular}
\par\smallskip
$\left.\begin{array}{l}\text{Sonnenferne}\\\text{Knoten}\end{array}\right\}\ \text{siderisch ruhend}\left\{\begin{array}{lr}
.\ .\ .\ .\ .\ .\ .\ .\ .&300^0\quad5'\quad4''\\
.\ .\ .\ .\ .\ .\ .\ .\ .&172\ 34\ 35
\end{array}\right.$
\par\smallskip
\begin{tabular}{l r}
Excentricität $\ .\ .\ .\ .\ .\ .\ .\ .\ .\ .\ .\ .\ .\ .\ .\ .$&$0.2591096$\\
Neigung $\ .\ .\ .\ .\ .\ .\ .\ .\ .\ .\ .\ .\ .\ .\ .\ .\ .\ .$&$35^0\quad0'\ 42''$
\end{tabular}}
\end{center}

Mit diesen Elementen verglich nun Dr.\ \textsc{Gauss} die ganze Reihe unserer \textit{Seeberger} Beobachtungen, und fand ihre Übereinstimmung, wie hier folgt.
\begin{center}
{\small\setlength{\tabcolsep}{5pt}
\begin{tabular}{r r|r|r|r|r}
\multicolumn{2}{c|}{}&&&\multicolumn{2}{c}{Unterschiede}\\
\multicolumn{2}{c|}{1802}&\shortstack{Berechnete\\AR.\\der Pallas}&\shortstack{Berechnete\\Abweichung\\der Pallas\\nördlich}&in AR.&in Decl.\\\hline
April&4&$183^0\ 44'\quad5''7$&$13^0\ 54'\ 55''4$&$-0''9$&$+3''4$\\
&5&$183\ 34\ 24.6$&$14\ 13\ 22.6$&$+0.9$&$-0.3$\\
&7&$183\ 15\ 39.6$&$14\ 49\quad1.3$&$+1.6$&$-0.7$\\
&8&$183\quad6\ 35.4$&$15\quad6\ 15.2$&$-2.4$&$+5.2$\\
&15&$182\ 10\ 17.9$&$16\ 54\ 35.7$&$+1.4$&$+4.9$\\
&18&$181\ 50\ 29.4$&$17\ 34\ 27.1$&$-1.2$&$\ .\ .\ .$\\
&19&$181\ 44\ 30.7$&$17\ 46\ 51.9$&$+5.4$&$-2.5$\\
&24&$181\ 19\ 35.8$&$18\ 42\ 27.9$&$-2.0$&$\ .\ .\ .$\\
&25&$181\ 15\ 38.3$&$18\ 52\ 18.9$&$+6.1$&$\ .\ .\ .$\\
&26&$181\ 12\quad2.0$&$19\quad1\ 44.9$&$+0.2$&$-4.5$\\
&27&$181\quad8\ 46.6$&$19\ 10\ 46.3$&$+1.0$&$-0.5$
\end{tabular}
\par\medskip
\begin{tabular}{r r|r|r|r|r}
\multicolumn{2}{c|}{}&&&\multicolumn{2}{c}{Unterschiede}\\
\multicolumn{2}{c|}{1802}&\shortstack{Berechnete\\AR.\\der Pallas}&\shortstack{Berechnete\\Abweichung\\der Pallas\\nördlich}&in AR.&in Decl.\\\hline
April&29&$181^0\quad3'\ 21''2$&$19^0\ 27'\ 36''2$&$+4''6$&$-7''5$\\
&30&$181\quad1\ 11.4$&$19\ 35\ 24.9$&$+1.5$&$-8.6$\\
Mai&1&$180\ 59\ 22.9$&$19\ 42\ 50.8$&$+0.4$&$+10.0$\\
&2&$180\ 57\ 56.7$&$19\ 49\ 53.2$&$+4.8$&$+1.6$\\
&3&$180\ 56\ 52.4$&$19\ 56\ 32.9$&$-3.2$&$-14.8$\\
&5&$180\ 55\ 49.6$&$20\quad8\ 45.6$&$-5.0$&$+6.9$\\
&6&$180\ 55\ 51.4$&$20\ 14\ 19.4$&$-2.7$&$\ .\ .\ .$\\
&7&$180\ 56\ 14.9$&$20\ 19\ 32.1$&$-8.9$&$-3.2$\\
&8&$180\ 57\quad0.4$&$20\ 24\ 23.7$&$-9.6$&$+3.6$\\
&11&$181\quad1\ 25.5$&$20\ 36\ 58.8$&$-16.2$&$-14.4$
\end{tabular}}
\end{center}
\SourceLeaf{220}%

Sehr leicht hätte sich Dr.\ \textsc{Gauss} an die letzten geraden Aufsteigungen noch etwas genauer anschliessen können; aber wegen der, bei anrückender Dämmerung, weniger zuverlässigen Declinationen würde er doch nicht sicher gewesen sein, ob er dadurch wirklich eine \textit{Verbesserung} der Bahn erhielte. Er will daher noch spätere Beobachtungen abwarten, ehe er eine neue Correction dieser Elemente vornehmen wird.

Dem sei nun wie ihm wolle, so scheinen doch alle bisherige Beobachtungen, Berechnungen und Vergleichungen den Planetismus der \textit{Pallas} zu begründen, und uns gegen alle Zweifel, Bedenklichkeiten und unbestimmte Hypothesen zu schützen. Denn, wie Dr.\ \textsc{Gauss} in einem Schreiben ganz gut bemerkt: „Einem scharfen Calcul kann man doch nichts entgegensetzen, als einen eben so scharfen Calcul, und nicht leere Vermuthungen und vage Raisonnemens. Allerdings werden zu einer leichten Wiederauffindung der Pallas im Jahr 1803 meine II. Elemente noch ansehnliche Verbesserungen bedürfen, und hoffentlich auch erhalten, wenn nur noch hinlängliche Beobachtungen dazu glücken; allein \textit{ansehnlich} werden diese Verbesserungen nur in Rücksicht der Bestimmung des Orts für 1803 sein, aber nicht in Ansehung des \textit{Planetismus}, oder \textit{Nicht-Planetismus} der \textit{Pallas}.“

Schon bei den früher gefundenen Elementen fiel die Bahn der \textit{Pallas}, da wo ihr aufsteigender Knoten auf der Ceres-Bahn ist, sehr hart an die Ceres-Bahn. Dies brachte Dr.\ \textsc{Olbers} auf die Vermuthung eines wirklichen Schnittes der \textit{Pallas}- und \textit{Ceres-Bahn}; dass beide sich etwa in einander, wie zwei Glieder einer Kette, schlingen. Nach einem Überschlage, welchen Dr.\ \textsc{Gauss} für die Abstände der \textit{Pallas} und \textit{Ceres} in der Knotenlinie ihrer Bahnen gemacht hat, fand er für $\DescendingNode$ den Abstand der \textit{Pallas} von der Sonne $2.86$, und der \textit{Ceres} $2.93$; beim $\AscendingNode$ sind diese Abstände weniger gleich. Die Reihe der Beobachtungen ist aber noch zu kurz, um die Möglichkeit, dass die Distanzen an dieser Stelle vielleicht genau gleich sind, läugnen zu können. Dr.\ \textsc{Gauss} will der Rechnung ganz ihren Gang lassen, und schlechterdings nichts hypothetisches einmischen. Es ist freilich noch zu früh, Hypothesen nachzuhängen, da wir statt derselben, wenn künftiges Jahr die Wiederauffindung der \textit{Pallas} glückt, bald Gewissheit bekommen werden.

Indessen, zu welchen Speculationen über die Entstehung und Geschichte unseres Planetensystems wird dieser sonderbare Weltkörper nicht Anlass geben? Welch’ eine ganz unerwartet grosse Neigung für einen Planeten? und wie ganz auffallend die Lage seiner Bahn gegen die Ceres-Bahn? Diese brachte Dr.\ \textsc{Olbers} auf den Gedanken, dass beide Planeten wol nur \textit{Rudera} eines Einzigen, durch den Stoss eines Cometen zertrümmerten sein könnten. Wer hätte so etwas jemals in unserem Planetensysteme erwartet, und zu welchen neuen, wichtigen und grossen Aufschlüssen kann und wird uns dieser Weltkörper nicht führen, welcher in mehr als einer Rücksicht zu den allerwichtigsten, jemals gemachten astronomischen Entdeckungen gehört.

Eben so merkwürdig findet Dr.\ \textsc{Gauss} die grosse Übereinstimmung der mittlern Bewegung bei diesen beiden Planeten. Bis jetzt vermag man nicht mit Gewissheit zu entscheiden, ob sie vielleicht \textit{genau} gleich sind. Kennen wir doch die eigentliche tägliche. Bewegung der \textit{Ceres} noch nicht auf eine Secunde genau. So klein aber auch die Massen der \textit{Ceres} und \textit{Pallas} sind, so muss es doch für den Unterschied ihrer mittlern Bewegungen eine Gränze geben, innerhalb welche er nicht fallen kann, ohne durch die wechselseitige Einwirkung ganz auf Null gebracht zu werden. Nun sind aber diese mittleren Bewegungen nach den II. Elementen nur wenig verschieden, und vielleicht kommen sie bei genauerer Kenntniss der Bahnen einander noch viel näher. So könnte es doch vielleicht sein, dass \textit{Ceres} und \textit{Pallas} einerlei Umlaufszeit hätten, immer nicht sehr weit von einander entfernt wären, ohne doch einander jemals von Anbeginn an zu nahe gewesen zu sein, oder künftig zu nahe kommen zu können, da beim $\AscendingNode$ \textit{Pallas}, beim $\DescendingNode$ \textit{Ceres}, wegen der so sehr verschiedenen Mittelpunktsgleichung allemal voraus wäre. Man wird alsdann auch nach einer langen Reihe von Jahren aus Beobachtungen der \textit{Ceres} die
\SourceLeaf{221}%
Masse der \textit{Pallas}, aus Beobachtungen der \textit{Pallas} die Masse der \textit{Ceres} sehr gut bestimmen können. Dr.\ \textsc{Gauss} findet nach einem Überschlage, dass 12 Secunden Unterschied in der mittleren täglichen Bewegung eine synodische Revolution von beinahe drei hundert Jahren gibt.

\begingroup\rightskip=0pt plus1em\relax
Zu welchen Untersuchungen wird ein solches Phänomen, einzig in seiner Art, nicht Veranlassung geben? Welche neue Bedürfnisse der Wissenschaft wird es nicht wecken? Denn nach allen unseren bisherigen Kenntnissen hätten wir uns solcher ausserordentlicher Erscheinungen im Weltsystem nicht versehen! Wie werden uns unsere bisherigen üblichen Rechnungs-Methoden verlassen und unzureichend sein! z.\ B.\ wenn wir nur die Störungen dieses neuen Weltkörpers berechnen sollen. Bisher hat man, nach den allein üblichen Methoden, die Excentricität und Neigung der Bahnen gleichsam als unendlich kleine Grössen betrachtet, und die Entwickelung nach den Potenzen derselben, mit Vernachlässigung der höhern, vorgenommen. Dies wird aber nicht mehr der Fall sein können bei zwei so starken Elementen der Pallas-Bahn, deren höhere Potenzen gewiss nicht vernachlässigt werden dürfen. Daher gründet vielleicht auch die \textit{Pallas} nothwendig eine neue Epoche in der physischen Astronomie. Nur was wir durch strengen Calcul wissen, ist ja das einzige, was wir über den Weltbau \textit{wirklich wissen}. Aber der Calculator bedarf Erfahrungen und Beobachtungen; nur durch diese stempelt er Hypothesen zur Wahrheit. Die \textit{Pallas} ist daher von allen Astronomen mit eben so grossem Fleisse als Sorgfalt beobachtet worden. Wir theilen unseren Lesern die ganze gesammelte Reihe dieser kostbaren Beobachtungen mit. Unsere Zeitschrift war bisher das Archiv aller Beobachtungen der \textit{Ceres}; sie soll es auch für die \textit{Pallas} sein, und wir werden keine Mühe sparen, diese Sammlung so vollständig als möglich zu machen, damit die Astronomen hier alles beisammen finden, was sie nur aus zerstreuten Blättern mühsam zusammen suchen müssten, und nur zu spät, vielleicht auch nie, finden würden.\par\endgroup

—\quad —\quad —

Um denjenigen Astronomen, welche mit grossen Aequatorial-Instrumenten versehen sind, das Aufsuchen der \textit{Pallas} zu erleichtern, hat Dr.\ \textsc{Gauss} nach seinen II. Elementen der Bahn nachstehende Ephemeride ihres Laufes berechnet. Ob es gleich etwas gewagt scheint, aus Elementen, welche sich nur auf 27tägige Beobachtungen gründen, den Ort des Planeten damit auf 2 Monate voraus bestimmen zu wollen: so kann es zwar nicht fehlen, dass diese Ephemeride gegen das Ende schon sehr merklich vom Himmel abweichen wird; da aber diese Abweichung wahrscheinlich nur einige Minuten betragen kann: so wird dies die Brauchbarkeit derselben zur bequemen Auffindung nicht hindern, zumal da man schon voraussehen kann, dass der Fehler in AR (und vermuthlich auch in der Declin.) negativ sein wird.
\Needspace{12\baselineskip}
\begin{center}
\textit{Ephemeride für die Pallas Olbersiana, für Mitternacht in Seeberg.}\par\medskip
{\small\setlength{\tabcolsep}{5pt}
\begin{tabular}{r r|r|r@{\qquad}r r|r|r}
\multicolumn{2}{c|}{1802}&\shortstack{AR.\\der Pallas}&\shortstack{Declinat.\\der Pallas}&\multicolumn{2}{c|}{1802}&\shortstack{AR.\\der Pallas}&\shortstack{Declinat.\\der Pallas}\\\hline
Mai&24&$181^0\ 57'$&$21^0\quad1'$ N.&Juni&14&$185^0\ 12'$&$20^0\ 17'$ N.\\
&27&$182\ 18$&$21\quad0$&&17&$185\ 48$&$20\quad5$\\
&30&$182\ 41$&$20\ 57$&&20&$186\ 27$&$19\ 52$\\
Juni&2&$183\quad6$&$20\ 52$&&23&$187\quad7$&$19\ 37$\\
&5&$183\ 34$&$20\ 46$&&26&$187\ 49$&$19\ 22$\\
&8&$184\quad5$&$20\ 38$&&29&$188\ 32$&$19\quad6$\\
&11&$184\ 37$&$20\ 28$&&&&
\end{tabular}}
\end{center}
\begin{flushright}[\textsc{von Zach}.]\end{flushright}
